\bookchapter{第二十四章：Root権限とBLロックの深掘り解説}

\chapterintro{この章は第二十章の内容を深く掘り下げるもので、スマートフォンの基盤メカニズムを真に理解したい、あるいはフラッシュを試みたいユーザー向けです。これらの操作にはリスクがあるため、慎重に行ってください。}

\booksection{一、Bootloaderの深い理解}

：ユーザーが旧バージョンのシステムをフラッシュして既知の脆弱性を悪用するのを防ぎます。

\booksubsection{データ保護}

\begin{itemize}
  \item \textbf{：新システムをフラッシュしても、元のアカウントパスワードがなければユーザーデータにアクセスできません。}二、BLロック解除の詳細フロー（Xiaomiを例に）
  \item \textbf{ロック解除権限の申請}：Xiaomiのロック解除ツールサイトで解除を申請。
  \item \textbf{アカウント連携}：スマートフォンの設定でXiaomiアカウントを連携。
\end{itemize}

\booksubsection{ハードウェアレベルの保存と検証}

BL ロックの状態（Locked/Unlocked）およびメーカーの公開鍵は、通常保護されたハードウェア領域（Qualcomm の devinfo パーティション、eFuse や RPMB パーティションなど）に保存され、ソフトウェアによる改ざんを防ぎます。起動のたびに Bootloader はこのセキュア領域から状態を読み取り、厳格な署名検証を実行するかどうかを決定します。

\booksubsection{デジタル署名検証}

メーカーはシステム更新やイメージを発表する際、自身の秘密鍵で署名します。デバイス起動時、Bootloader は内蔵の公開鍵でシステムイメージ（boot.img、recovery.img など）を検証します。その仕組みは：

\begin{enumerate}
  \item システムイメージのハッシュ値を再計算。
  \item 公開鍵でイメージに付属の署名を復号し、ハッシュ値を取得。
  \item 2つのハッシュ値を比較し、一致すれば検証合格。これによりシステムイメージが第三者に改ざんされていないことを保証します。
\end{enumerate}

\booksubsection{セキュアブートチェーン}

BL ロックはセキュアブートチェーンの第一段階です。検証合格後、次段階のブートプログラム（カーネルなど）をロードして実行し、システム起動環境の信頼性を段階的に保証します。

\booksubsection{ロック vs 解除：BL ロックの両面影響}

BL ロックの利害をより直感的に理解するため、5つのコア次元での違いを対比表で示します：

\begin{center}
\begin{tabularx}{\linewidth}{|X|X|X|}
\hline
\textbf{核心機能} & \textbf{ロック状態（デフォルト）} & \textbf{解除後の状態} \tabularnewline
\hline
システムセキュリティ & 未署名の非公式システムのフラッシュを防止し、システム完全性と低層セキュリティを保護\newline 高セキュリティ：公式署名イメージのみロード、悪意あるファームウェア改ざんを効果的に防御 & セキュリティ低下：サードパーティ ROM のフラッシュを許可、システム完全性検証メカニズムが無効化、マルウェアが Root 権限を取得するとシステムを深く改ざん可能 \tabularnewline
システム安定性 & メーカーがテスト·最適化した公式システムの実行を保証し、安定性と快適さを確保\newline 安定：公式ファームウェア使用、機能完備、OTA 更新が信頼できる & リスク増大：サードパーティ ROM に互換性問題がある可能性、システムクラッシュ、文鎮化、ハードウェア機能異常（カメラ、急速充電など）、OTA 更新失敗の可能性 \tabularnewline
保証権益 & メーカーは BL ロックで安定した公式システムを提供し、保証サービスを提供\newline 保証有効：公式無料保証サービスを享受（非意図的損傷） & 保証無効：ほとんどのメーカーは BL 解除を能動的な保証放棄とみなします。再ロックしても履歴解除記録が検出されれば保証拒否される可能性 \tabularnewline
データセキュリティ & 解除プロセスは通常データの強制消去を行い、旧データが新システム環境で悪用されるのを防止\newline データ安全：日常使用中、BL ロックがマルウェアのシステムパーティション変更能力を制限 & データリスク：解除時に全ユーザーデータが強制消去；解除後に正しく暗号化しないとデータ保存が不安定；紛失後にフラッシュでデータを抽出されやすい \tabularnewline
機能制限 & セキュア環境に依存する機能（決済、銀行 App など）の正常動作を保証\newline 機能完備：すべての公式機能が正常に使用可能 & 機能制限：一部セキュリティ敏感 App（Google Pay、銀行 App など）が BL 解除状態を検出して実行を拒否する可能性；DRM サービスがダウングレードされ、高画質動画再生に影響 \tabularnewline
\hline
\end{tabularx}
\end{center}

\booksection{二、BLロック解除の詳細フロー（Xiaomiを例に）}

\begin{enumerate}
  \item \textbf{：アカウントランクに応じ7-30日の待機が必要。}解除操作
  \item \textbf{：Xiaomiロック解除ツールを使用し、パソコンに接続して解除。}三、Rootの複数の方案
  \item \textbf{Magisk}：現代的なRoot方案、Systemless設計、Root状態を隠蔽可能。
  \item \textbf{Kernelsu}：カーネルベースのRoot方案、より低層、隠蔽性が高い。
\end{enumerate}

\booksection{三、Rootの複数の方案}

\begin{itemize}
  \item \textbf{：古参のRoot方案、現在は更新停止。}四、フラッシュ入門
  \item \textbf{フラッシュとはスマートフォンに新しいOSをインストールすることです。一般的なフラッシュの種類：}公式ROMのフラッシュ
  \item \textbf{：公式システムのインストール、文鎮化からの復旧やダウングレードに使用。}サードパーティRecoveryのフラッシュ
\end{itemize}

\booksection{四、フラッシュ入門}

サードパーティROMのフラッシュ

\begin{itemize}
  \item \textbf{：LineageOS、Pixel Experienceなど、異なるシステムを体験。}ROM焼きとは、スマートフォンに新しいオペレーティングシステムをインストールすることです。一般的なROM焼きの種類：
  \item \textbf{公式ROM焼き}：公式システムをインストールし、文鎮復旧やダウングレードに使用。
  \item \textbf{サードパーティRecovery焼き}：TWRPなど、バックアップ復元やサードパーティROMの書き込みに使用。
\end{itemize}

\booksection{五、Samsung Knox：BL 解除の不可逆な代償}

他の Android メーカーと異なり、Samsung スマホの BL 解除メカニズムには独特の「ヒューズ」設計があります：\textbf{Samsung スマホにおいて、解除は Knox ヒューズ作動をトリガーし、一部セキュリティ機能（Samsung Pay など）を永久に喪失；保護データをすべて消去；保証に影響}。具体的には：

\begin{itemize}
  \item \textbf{ヒューズは使い切り}: スマホ内部に eFuse ヒューズがあり、BL 解除を検出すると永久に焼き切れ、復元できません。
  \item \textbf{影響を受ける機能}: Samsung Pay（交通系/銀行カード）、Knox セキュアフォルダー、企業 MDM 管理制御、一部健康保険系 App がすべて無効化されます。
  \item \textbf{データを強制消去}：ヒューズがトリガーされるとスマホは強制的に初期化され、すべてのローカルデータが消去されます（クラウドデータは影響を受けません）。
  \item \textbf{保証への影響}：Knoxヒューズトリガー後は保証サービスの対象外となることが公式に明示されています。
\end{itemize}

\begin{tcbwarn}{行動する前に慎重に}
SamsungのKnoxヒューズは\textbf{不可逆的}で、一度トリガーされると元の工場状態には戻せません。Root権限が必要な機能を使いたいだけなら、Root不要の代替手段があるかまず調べてください。

\end{tcbwarn}

\textbf{サードパーティROM焼き}\par

：LineageOS、Pixel Experienceなど、異なるシステムを体験。

\begin{itemize}
  \item 重要な警告
  \item ROM焼きにはリスクがあり、操作には慎重さが必要です。次のような問題が発生する可能性があります：
  \item スマートフォンが文鎮化（起動不可）
  \item 保証の喪失
\end{itemize}

データの消失

\begin{tcbinfo}{決済機能の異常}
\begin{itemize}
  \item \textbf{ROM焼きの前に必ず完全なバックアップを行い、チュートリアルを注意深く読み、機種とバージョンの一致を確認してください。}おすすめリソース
  \item \textbf{XDA Developers}：世界最大のAndroid開発者コミュニティ、豊富なROM焼きチュートリアルとROMリソースがあります。
  \item \textbf{Coolapk}：中国国内で有名なデジタルコミュニティ、中国語のチュートリアルや議論が豊富です。
\end{itemize}

\end{tcbinfo}

\booksection{六、よくある質問と回答（Q\& A）}

\textbf{Q1：BLロック解除後に元に戻せますか？}\par

\textbf{A：}再ロック可能です。通常、公式の完全なフラッシュパッケージを書き込み、特定のフラッシュスクリプト（Xiaomiの\texttt{flash\_all\_lock.\allowbreak bat}）を使って再ロックします。ただし、再ロック後も、解除によってトリガーされたセキュリティヒューズ（SamsungのKnoxなど）は元に戻らない場合があります。

\textbf{Q2：ロック解除後にRoot化するには？}\par

\textbf{A：}BLロック解除は最初のステップにすぎません。Root化には通常以下が必要です：

\begin{enumerate}
  \item カスタムRecovery（TWRPなど）のフラッシュ。
  \item Recovery経由でMagiskなどのRoot管理ツールをフラッシュし、bootイメージをパッチしてRoot権限を取得。
\end{enumerate}
